% Generated draft: requires mathematical and source review.
\documentclass[11pt]{article}
\usepackage[margin=25mm]{geometry}
\usepackage{fontspec}
\setmainfont{Latin Modern Roman}
\usepackage{amsmath,amssymb,mathtools,unicode-math}
\setmathfont{Latin Modern Math}
\usepackage{enumitem,needspace,xurl,hyperref}
\setlength{\parindent}{0pt}
\setlength{\parskip}{5pt}
\setlength{\emergencystretch}{3em}
\newcommand{\N}{\mathbb{N}}
\newcommand{\Z}{\mathbb{Z}}
\newcommand{\Q}{\mathbb{Q}}
\newcommand{\R}{\mathbb{R}}
\newcommand{\C}{\mathbb{C}}
\newcommand{\F}{\mathbb{F}}
\newcommand{\A}{\mathbb{A}}
\newcommand{\PP}{\mathbb{P}}
\newcommand{\E}{\mathbb{E}}
\newcommand{\eps}{\varepsilon}
\newcommand{\dd}{\,\mathrm{d}}
\begin{document}
\section*{UnsolvedMath ID: 30006172}
\textbf{Source problem code:} catalogue entry
\subsection{Definitions and mathematical statement}
Let \(\omega_1\) be the least uncountable ordinal and \(\omega_2\) the next initial ordinal. Set
\[
X={}^{\omega_1}\omega_1=\{f\mid f:\omega_1\to\omega_1\}.
\]
The supplied record refers to a particular ``function-space measure'' \(\nu\) on \(X\) and its associated elementary embedding \(j_\nu\). It does not define this measure, specify the background axioms, or give the domain, codomain, or construction of the embedding. These are unresolved prerequisites; the statement below concerns the particular objects intended by the source.

\textbf{Requested computation.} For that intended measure and embedding, determine the ordinal
\[
j_\nu(\omega_2).
\]
The reconstructed statement additionally asks for the values of \(j_\nu\) on ``relevant cardinals,'' but supplies no class of cardinals or precise quantifier for this extension.

\textbf{Supplied contextual bound.} The extracted original passage attributes to Chan a strict upper bound for \(j_\nu(\omega_1)\). Its right-hand side is rendered as ``\(\delta\) 13'' in the supplied extraction. This may denote \(\delta^1_3\), but the indices and the intended invariant are not established by the supplied record. Consequently, neither that interpretation nor a definition of the bound is asserted here. The attribution and strictness of the reported inequality are retained as source information, not independently verified mathematical facts.
\subsection{Short English statement}
Determine \(j_\nu(\omega_2)\) for the particular function-space measure on \({}^{\omega_1}\omega_1\) intended by the source. The measure, background axioms, and associated embedding require clarification. The broader request about other cardinals and the notation in a bound attributed to Chan also remain unresolved.
\subsection{Sources}
\begin{itemize}
\item Local UnsolvedMath registry: supplied source record, permanent ID \(30006172\), exact source problem code \texttt{catalogue entry}, title ``Function-Space Measure Embedding Values.'' The record belongs to the Oberwolfach Reports Open Problems collection and gives proposed year 2025; no proposer is identified.
\item Supplied report citation: \emph{Set Theory} (2025), Oberwolfach Reports, DOI 10.4171/owr/2025/2, pp.\ 75--140. \url{https://doi.org/10.4171/owr/2025/2}.
\item Supplied literature source: ``Oberwolfach Report 2/2025.'' \url{https://content.ems.press/assets/public/full-texts/serials/owr/22/1/14299082/online/10.4171-owr-2025-2.pdf}.
\item Literature assessment dated exactly \texttt{2026-08-22}: the supplied curation labels the problem open and says the report requests the next embedding value while reporting only a strict upper bound for a lower value. This is supplied dated curation, not an independently verified assessment of the literature or current status.
\item The registry labels the statement \texttt{corrected\_verified} and says it was reconstructed and checked against the official report. Those are supplied provenance claims; no independent source verification is asserted here. The extracted passage attributes the lower-value bound to Chan, without a full name or separate bibliographic reference.
\end{itemize}
\end{document}
